\begin{table}[tp]\centering\footnotesize\setlength{\tabcolsep}{3pt}
\caption{Implied Trade-Off: The Phillips Multiplier}\label{tab:tradeoff}
\resizebox{\ifdim\width>\textwidth\textwidth\else\width\fi}{!}{%
\begin{tabular}{lcccc}\toprule
 & \multicolumn{2}{c}{24 months} & \multicolumn{2}{c}{36 months}\\ \cmidrule(lr){2-3}\cmidrule(lr){4-5}
Yield and FCI, measured as & Estimate & 90\% set & Estimate & 90\% set\\ \midrule
Nominal & $-$1.21 & undefined; 68\%: [$-$3.29, $-$0.30] & $-$1.03 & undefined; 68\%: [$-$2.96, $-$0.31]\\
Real & $-$1.16 & [$-$5.51, $-$0.25] & $-$0.90 & [$-$7.04, $-$0.07]\\
\textbf{Neutral rate, 1996 start (baseline)} & \textbf{$-$0.95} & \textbf{[$-$1.99, $-$0.32]} & \textbf{$-$0.81} & \textbf{[$-$2.12, $-$0.18]}\\
Neutral rate, 1973 start & $-$1.49 & [$-$2.92, $-$0.72] & $-$1.64 & [$-$5.39, $-$0.72]\\
\bottomrule\end{tabular}}
\begin{flushleft}\footnotesize\textit{Notes:} $\kappa_H$ is the core PCE price level at month $H$ (percent) per percentage-point year of cumulative unemployment-gap response \citep{barnichon2021}; its inverse in absolute value is the sacrifice ratio. Estimate: ratio of the point responses. 90\% set: Anderson--Rubin type, all values for which the 90\% bootstrap interval of $\tilde P_H-\kappa\tilde U_H$ contains zero. The ratio is undefined when the 90\% interval of the cumulative unemployment response includes zero, because the data then cannot rule out unemployment responses of zero or either sign; the 68\% set is then shown where it exists. 12 lags. All cases in Appendix Table~\ref{tab:tradefull}.\end{flushleft}
\end{table}
